%% Template: ats-resume-classic
%% A single-column, icon-free, table-free U.S. resume modeled directly on
%% a real compact classic resume layout: one-line contact header, Summary
%% -> Professional Experience -> Education -> Skills and Certifications
%% (skills/certs as terse comma lists at the end, not a verbose narrative
%% "Core Competencies" block). Built to genuinely fit 1 page for a
%% candidate with 3-4 real roles, not just "allow up to 2".
%%
%% Compile with: lualatex (xelatex also works unchanged)
%% Compile command: lualatex -interaction=nonstopmode <file>.tex

\documentclass[10.5pt,letterpaper]{article}

\usepackage[margin=0.6in]{geometry}
% fontspec + TeX Gyre Heros: a free, metric-compatible Helvetica clone.
% Helvetica/Arial-style sans-serif is one of the most commonly
% recommended ATS-safe resume fonts (clean, highly legible at small
% sizes, universally familiar to recruiters), and TeX Gyre Heros ships
% as standard OpenType with every TeX Live/MiKTeX install -- no bundled
% font files needed. Verified other common resume-font picks (Carlito/
% Calibri clone, Lato, Open Sans, PT Sans) are NOT available on a
% stock MiKTeX install without adding font packages; only tex-gyre
% fonts (Heros/Termes/etc.) are guaranteed present out of the box.
% Needs no xltxtra/xunicode (XeTeX-era packages fontspec alone has long
% since made unnecessary), so this carries none of the fontspec/
% xunicode MiKTeX auto-installer fragility other custom-font templates
% in this repo can hit.
\usepackage{fontspec}
\setmainfont{TeX Gyre Heros}
\usepackage[hidelinks]{hyperref}
\usepackage{titlesec}
\usepackage{enumitem}
\usepackage{needspace}

% Fully monochrome -- no xcolor, no accent color, matches a plain
% printed-Word-doc look. ALL CAPS bold section labels, no rule line
% underneath (deliberately -- the reference resume this is modeled on
% uses whitespace, not a rule, to separate sections, which reads as
% cleaner at 10.5pt and saves vertical space toward the 1-page target).
\titleformat{\section}{\bfseries\MakeUppercase}{}{0em}{}
\titlespacing{\section}{0pt}{6pt}{2pt}

% Plain ASCII "-" (a single hyphen, NOT "--") as the bullet marker.
% "--" is a real LaTeX gotcha: two consecutive hyphens ligature into an
% en dash (-) at typeset time, not a plain hyphen -- an earlier version
% of this template used "label=--" here by mistake, which is why
% bullets and headers rendered with dash-like characters instead of
% plain hyphens. A single "-" never ligatures. This also sidesteps a
% separate risk with \textbullet: under some font/engine combinations
% it can lack a ToUnicode mapping and extract as an unreadable
% replacement character in pdftotext. Tight spacing throughout
% (itemsep/topsep/parsep all near-zero) is what makes the 1-page target
% achievable with a full 3-4 job history.
\setlist[itemize]{leftmargin=1.05em, itemsep=0pt, topsep=1pt, parsep=0pt, label=-}

\setlength{\parindent}{0pt}
\pagenumbering{gobble}

% ------------------------------------------------------------------
% \cvheader{Title}{Company, Location}{Dates}
% Plain-text two-line block -- no tabular/tabularx, and deliberately no
% \hfill right-alignment either: \hfill-based date columns were tested
% and found to misattribute a date to the wrong entry under
% `pdftotext -layout` (looks fine visually, breaks silently in the text
% layer). Dates go inline right after the title instead, which is the
% only ordering pdftotext extracts unambiguously. Separator is a pipe
% ("|", matching the header contact line), not a dash -- avoids any
% dash-like character in the title/date line entirely.
% ------------------------------------------------------------------
\newcommand{\cvheader}[3]{%
  \noindent\textbf{#1} \textbar\ \textit{#3}\\
  \textit{#2}
}

\begin{document}

% ============================================================
%     HEADER
% ============================================================
\begin{center}
  {\Large\bfseries [YOUR_NAME]}\\[1pt]
  [YOUR_CITY], [YOUR_STATE] \textbar\ [YOUR_PHONE] \textbar\ \href{mailto:[YOUR_EMAIL]}{[YOUR_EMAIL]} \textbar\ \href{[YOUR_LINKEDIN_URL]}{[YOUR_LINKEDIN_DISPLAY]}
\end{center}
\vspace{2pt}

% Show the LinkedIn (and any other profile) URL as visible, literal
% text in the header above -- not hidden behind a generic "LinkedIn"
% link label. A printed copy or a non-clickable PDF viewer should still
% show the actual address, and some ATS keyword scans look for the
% literal domain/handle text.

% ============================================================
%     SUMMARY
% ============================================================
\section{Summary}
[2-4 line summary: years of experience, the 2-3 tools/skills most relevant to the target posting, and one quantified achievement. Keep this tight -- every line here is a line the Experience section doesn't get.]

% ============================================================
%     PROFESSIONAL EXPERIENCE
% ============================================================
\section{Experience}

\needspace{8\baselineskip}
\cvheader{[Job Title]}{[Company], [City, State]}{[Start]-[End]}
\begin{itemize}
  \item [Measurable achievement or responsibility -- action verb + number where possible]
  \item [Measurable achievement or responsibility]
  \item [Measurable achievement or responsibility]
  \item [Measurable achievement or responsibility]
\end{itemize}

\vspace{3pt}
\needspace{7\baselineskip}
\cvheader{[Job Title]}{[Company], [City, State]}{[Start]-[End]}
\begin{itemize}
  \item [Measurable achievement or responsibility]
  \item [Measurable achievement or responsibility]
  \item [Measurable achievement or responsibility]
\end{itemize}

\vspace{3pt}
\needspace{6\baselineskip}
\cvheader{[Job Title]}{[Company], [City, State]}{[Start]-[End]}
\begin{itemize}
  \item [Measurable achievement or responsibility]
  \item [Measurable achievement or responsibility]
\end{itemize}

\vspace{3pt}
\needspace{5\baselineskip}
\cvheader{[Job Title]}{[Company], [City, State]}{[Start]-[End]}
\begin{itemize}
  \item [Measurable achievement or responsibility]
\end{itemize}

% ============================================================
%     EDUCATION
% ============================================================
\section{Education}
\needspace{3\baselineskip}
\cvheader{[Degree, Field]}{[Institution], [City, State]}{[Start]-[End]}

% ============================================================
%     SKILLS AND CERTIFICATIONS
% ============================================================
% Deliberately terse comma lists, not full-sentence bullets -- this is
% the section that makes the 1-page target achievable versus a longer
% narrative "Core Competencies" block. Add/remove category lines to fit
% the candidate's actual skill set; do not pad with generic terms.
\section{Skills and Certifications}
\textbf{[Category 1]:} [comma-separated items] \\
\textbf{[Category 2]:} [comma-separated items] \\
\textbf{[Category 3]:} [comma-separated items] \\
\textbf{Certifications:} [Cert 1], [Cert 2], [Cert 3]

\end{document}
